\section{Explicit and recursive sequences}

A sequence is a function on natural-number inputs. To define such a function, we must say how its terms are obtained. There are two common ways to do this: an explicit formula and a recursive definition.

Both are useful, but they answer different practical questions.

\subsection*{Explicit formulas}

An explicit formula gives \(a_n\) directly in terms of \(n\). If we want the tenth term, we substitute \(n=10\). If we want the hundredth term, we substitute \(n=100\).

\begin{definitionbox}
An explicit formula for a sequence gives the term \(a_n\) directly as an expression depending on the index \(n\).
\end{definitionbox}

\begin{examplebox}
Let
\[
a_n=3n-2,
\qquad n\ge1.
\]
Then
\[
a_1=3\cdot1-2=1,
\]
\[
a_2=3\cdot2-2=4,
\]
\[
a_3=3\cdot3-2=7.
\]
The sequence begins
\[
1,\;4,\;7,\;10,\;\ldots
\]
If we want the tenth term, we compute
\[
a_{10}=3\cdot10-2=28.
\]
\end{examplebox}

The advantage of an explicit formula is direct access. We can compute a distant term without computing all previous terms.

\subsection*{Recursive definitions}

A recursive definition describes how to obtain the next term from earlier terms. Such a definition needs a starting value. Without a starting value, the sequence is not determined.

\begin{definitionbox}
A recursive definition of a sequence gives one or more initial terms and a rule for computing later terms from earlier terms.
\end{definitionbox}

\begin{examplebox}
Define a sequence by
\[
a_1=5,
\qquad
a_{n+1}=a_n+2.
\]
Then
\[
a_1=5,
\]
\[
a_2=a_1+2=7,
\]
\[
a_3=a_2+2=9,
\]
\[
a_4=a_3+2=11.
\]
The sequence begins
\[
5,\;7,\;9,\;11,\;\ldots
\]
\end{examplebox}

The recursive rule says how to move from one term to the next. It does not immediately give \(a_{100}\) unless we either compute many previous terms or find an explicit formula.

\subsection*{Why the initial term matters}

The recurrence relation alone may not define one sequence. It may define a whole family of possible sequences, depending on the starting value.

\begin{examplebox}
The rule
\[
a_{n+1}=a_n+2
\]
does not determine a unique sequence by itself.

If \(a_1=1\), the sequence begins
\[
1,\;3,\;5,\;7,\;\ldots
\]
If \(a_1=10\), the sequence begins
\[
10,\;12,\;14,\;16,\;\ldots
\]
The recursive rule is the same, but the sequences are different because the starting value is different.
\end{examplebox}

This is similar to describing motion. A rule for how to move is not enough if we do not know where we start.

\subsection*{Geometric-type recursion}

Some recursions multiply instead of add.

\begin{examplebox}
Define
\[
a_1=3,
\qquad
a_{n+1}=2a_n.
\]
Then
\[
a_1=3,
\]
\[
a_2=2a_1=6,
\]
\[
a_3=2a_2=12,
\]
\[
a_4=2a_3=24.
\]
The sequence begins
\[
3,\;6,\;12,\;24,\;\ldots
\]
Each term is twice the previous one.
\end{examplebox}

The rule is simple, but the meaning is important. The next value depends on the current value. The sequence is built step by step.

\subsection*{Reading recursive notation}

The expression \(a_{n+1}\) means the term after \(a_n\). If \(n=1\), then \(a_{n+1}=a_2\). If \(n=2\), then \(a_{n+1}=a_3\).

A recurrence relation such as
\[
a_{n+1}=a_n+2
\]
should be read as a rule that works repeatedly:
\[
a_2=a_1+2,
\qquad
a_3=a_2+2,
\qquad
a_4=a_3+2,
\]
and so on.

\begin{remarkbox}
Recursive notation is compact because one line contains infinitely many instructions. To understand it, expand the first few cases and read what the rule is doing.
\end{remarkbox}

\begin{summarybox}
When reading a sequence definition, ask:
\begin{itemize}
\item Is the sequence given explicitly or recursively?
\item If it is explicit, can I compute \(a_n\) directly from \(n\)?
\item If it is recursive, what is the starting value?
\item How is the next term obtained from earlier terms?
\item What are the first few terms?
\item Does the notation \(a_{n+1}\) refer to the next term after \(a_n\)?
\end{itemize}
\end{summarybox}

Explicit and recursive descriptions are two languages for sequences. Each one reveals a different aspect of the same kind of object: a function defined on natural-number inputs.